\documentclass[11pt,a4paper]{article}

% =============================
% Packages
% =============================
\usepackage[margin=0.85in]{geometry}
\usepackage{setspace}
\usepackage{graphicx}
\usepackage{booktabs}
\usepackage{longtable}
\usepackage{tabularx}
\usepackage{array}
\usepackage{float}
\usepackage{xcolor}
\usepackage{fancyhdr}
\usepackage{titlesec}
\usepackage{caption}
\usepackage{hyperref}
\usepackage{amsmath}
\usepackage{enumitem}
\usepackage{ragged2e}
\usepackage{makecell}
\usepackage{siunitx}

% =============================
% Document formatting
% =============================
\onehalfspacing
\hypersetup{colorlinks=true, linkcolor=black, urlcolor=blue}
\setlength{\parskip}{0.55em}
\setlength{\parindent}{0pt}
\renewcommand{\arraystretch}{1.18}
\newcolumntype{Y}{>{\raggedright\arraybackslash}X}
\sisetup{detect-all}

% =============================
% Section formatting
% =============================
\titleformat{\section}{\Large\bfseries}{\thesection}{0.7em}{}
\titleformat{\subsection}{\large\bfseries}{\thesubsection}{0.7em}{}

% =============================
% Header/footer for subsequent pages
% =============================
\pagestyle{fancy}
\fancyhf{}
\fancyhead[L]{\today}
\fancyhead[R]{Generated by LAMBDA: \url{https://lambda.com.ai}}
\cfoot{\thepage}
\renewcommand{\headrulewidth}{0.4pt}

% =============================
% First page custom header
% =============================
\fancypagestyle{firstpage}{
    \fancyhf{}
    \fancyhead[L]{%
        \raisebox{-0.5\height}{
\includegraphics[height=1.5cm]{./data/workspaces/17/.skills/res/logo.png}}
        \\[0.8em]
        \rule{\textwidth}{0.6pt}
    }
    \cfoot{\thepage}
    \renewcommand{\headrulewidth}{0pt}
}

\title{Pulmonary Hypertension Care Quality in Great Britain: \\A Data Analysis of the NAPH 10\textsuperscript{th} Annual Report}
\author{LAMBDA}
\date{10 August 2026}

\begin{document}
\thispagestyle{firstpage}

\begin{center}
\vspace*{2.5cm}
{\LARGE \textbf{Pulmonary Hypertension Care Quality in Great Britain}} \\[0.5em]
{\Large A Data Analysis of the NAPH 10\textsuperscript{th} Annual Report} \\[0.5em]
{\large LAMBDA} \\[0.5em]
{\large 10 August 2026}
\end{center}

\begin{abstract}
This report analyses data from the National Audit of Pulmonary Hypertension (NAPH) 10\textsuperscript{th} Annual Report (2018--19), covering eight specialist pulmonary hypertension centres across England, Scotland, and Wales. The dataset encompasses ten audit standards tracked over four years, active patient counts from 2010 to 2019, and detailed longitudinal survival data for patients diagnosed with idiopathic pulmonary arterial hypertension (IPAH). Across all measures, the national standard performance is strong: nine out of ten process quality standards were met at the national level in 2018--19, representing steady improvement since 2015--16. However, one critical area stands apart -- waiting times for pulmonary endarterectomy surgery fall dramatically short of targets, with only 11\% of cases meeting the four-month threshold against a 90\% target. Survival analysis reveals stark gradients by sex, age, and clinical severity, with five-year cumulative survival ranging from 80\% for young women without comorbidities to just 25\% for patients presenting with IPAH plus comorbidities. This report presents these findings with evidence from the audited data and discusses their implications for service planning and patient care pathways.
\end{abstract}

\newpage

\tableofcontents
\newpage

\section{How Well Do PH Centres Deliver on Agreed Standards?}

The National Audit of Pulmonary Hypertension measures clinical care against professionally agreed standards published by the Pulmonary Hypertension Outcomes Group (PHOG). In the 2018--19 audit year, nine of the ten adult-focused standards were met at the national level. This represents continued progress over the four-year trend period, with most standards showing either stable high performance or gradual improvement.

\begin{figure}[H]
  \centering
  \includegraphics[width=\textwidth]{./data/workspaces/17/chart_standards_dashboard.png}
  \caption{National standards performance at 2018--19, showing each standard's national achievement rate and the proportion of the seven adult centres that met the target. Nine of ten standards were met nationally; only Standard 13 (PEA wait times) falls significantly below its target. Green bars indicate the national figure meets or exceeds the target; red indicates failure to meet it.}
\end{figure}

Five standards achieved universal compliance across all seven adult centres in 2018--19: diagnosis recording (Standard 9), first-line PDE5 inhibitor use (Standard 10), quality of life documentation (Standard 11), annual consultations (Standard 12), and timely assessment for new patients (Standard 4, where all centres exceeded the 50\% floor target). Two additional standards -- timely diagnosis (Standard 3, 99\%) and cardiac catheterization before treatment (Standard 5, 97\%) -- were met nationally but one adult centre fell slightly short for each. 

The two areas requiring attention are drug therapy initiation within 12 weeks (Standard 8, 81\% nationally) and pulmonary endarterectomy waiting times (Standard 13, 11\%). While Standard 8's 81\% rate barely misses its 80\% target, only four of seven centres actually achieved this benchmark. Standard 13 remains the outlier: even though the national figure of 11\% fluctuates around 11--18\% across recent years, it consistently fails to approach the 90\% target by more than 70 percentage points.

The local action plans submitted by individual centres confirm that both of these shortcomings are recognised and actively addressed. For example, Golden Jubilee reduced its Standard 8 performance from 84\% in 2016--17 to 61\% in 2018--19 -- a concerning downward trajectory -- and explicitly states it is ``tackling this problem by increasing throughput for diagnostic right heart catheterisation.'' Meanwhile, Royal Free improved its Standard 8 rate from 71\% to 87\% between 2016--17 and 2018--19 by introducing targeted monitoring.

\begin{figure}[H]
  \centering
  \includegraphics[width=\textwidth]{./data/workspaces/17/chart_national_trends.png}
  \caption{Performance trends across nine nationally reported standards from 2015--16 through 2018--19. Most indicators show upward trajectories toward their targets until stabilising near or above the threshold. Standard 13 (PEA wait times, red line) deviates sharply from every other pattern, never coming close to its 90\% target despite minor fluctuations between 11\% and 18\%. Standard 8 (Drug Therapy within 12 Weeks, purple line) plateaus at approximately 81\%.}
\end{figure}

The upward trend in Standard 4 -- seen within 30 days -- is particularly notable. It rose from 47\% in 2015--16 to 63\% in 2018--19, an improvement of 16 percentage points over four years. Newcastle demonstrated the steepest gain, climbing from 19\% to 78\% in the same period, while Imperial College was the sole centre whose rate declined slightly (from 71\% to 62\%). The fact that all eight centres (including the paediatric site) achieved at least 50\% compliance in 2018--19 -- the first time this universal threshold was reached -- suggests genuine system-wide change rather than statistical noise.

\section{Patient Numbers Are Growing Substantially Across All Centres}

The total number of active patients managed by specialist PH centres increased from 4,341 in 2010 to 8,351 in 2019 -- an increase of 92\% over the decade. When the dataset is disaggregated by centre, the growth pattern holds broadly but with important asymmetries that may affect service planning.

\begin{figure}[H]
  \centering
  \includegraphics[width=\textwidth]{./data/workspaces/17/chart_patient_growth.png}
  \caption{Active patient population growth at adult PH centres from 2010 to 2019. Sheffield manages the largest caseload, while other centres follow similar growth trajectories. The national total nearly doubled from 4,341 to 8,351 over the decade. Notably, Sheffield's share of the national total has grown from approximately 31\% in 2010 to roughly 26\% in 2019 -- still the largest single centre.}
\end{figure}

Sheffield clearly carries the heaviest volume, managing 2,193 patients in 2018--19 -- almost three times as many as the next largest centre (Royal Papworth at 1,251). This concentration raises questions about whether Sheffield's workload per clinician and nurse might differ meaningfully from centres with smaller lists, although Sheffield did not participate in some standards reporting (implying it is likely excluded from certain comparisons due to data structure).

Great Ormond Street handles the smallest paediatric caseload at 292 active patients. Among adult centres, Newcastle consistently had the lowest volume, yet still comfortably exceeded the 300-patient minimum recommended by European Society of Cardiology / European Respiratory Society guidelines. All seven adult centres met Standard 2 (patient volume target of $\geq$300 PAH/CTEPH patients annually) throughout the audit period.

The growth curve does not appear to be levelling off. The absolute increase in managed patients accelerated between 2015 and 2019 (adding 1,401 patients over four years, compared with 916 added between 2010 and 2014). Commissioning bodies should plan accordingly, because the existing service infrastructure may come under pressure if referral patterns continue unabated.

\section{Waiting Times for Pulmonary Endarterectomy Are the Critical Weakness}

National Standard 13 requires that 90\% of patients who undergo pulmonary endarterectomy (PEA) -- the potentially curative surgery for chronic thromboembolic pulmonary hypertension (CTEPH) -- wait no longer than four months from diagnosis to operation. In 2018--19, only 11\% of cases achieved this benchmark nationally, making it the single worst-performing standard by an enormous margin.

\begin{figure}[H]
  \centering
  \includegraphics[width=\textwidth]{./data/workspaces/17/chart_pea_wait_times.png}
  \caption{Proportion of CTEPH patients meeting the four-month wait time target at each adult PH centre. Only Golden Jubilee and Royal Papworth consistently achieve at least partial compliance (23\% in each of three consecutive years), while several centres record 0\%. The dashed red line shows the 90\% target. The median interval from diagnosis to PEA surgery has improved since 2016 according to the main report, but the majority of patients still experience unacceptable delays.}
\end{figure}

Of the three adult centres that submit PEA data (Golden Jubilee, Imperial College, Royal Papworth, and Sheffield), only Golden Jubilee and Royal Papworth record any cases meeting the target. Both maintain a flat 23\% compliance rate across all three years -- meaning they do not appear to have worsened, but they also have not made meaningful progress. Newcastle, Royal Brompton and Harefield, and Royal Free all recorded 0\% compliance in 2018--19. These centres manage CTEPH patients who would require referral to the national surgical centre at Royal Papworth Hospital for endarterectomy, creating a complex multi-institutional pathway that current arrangements cannot sustain at scale.

Local action plans acknowledge the difficulty. Sheffield states the issue requires ``joint effort'' between NHS England and PH centres. Several centres refer directly to the Royal Papworth response, indicating that the bottleneck lies upstream of individual centre control -- specifically in surgical capacity constraints and pre-operative work-up delays. As the lead commissioner Kathy Blacker notes in the report foreword, commissioners and PH centres are conducting pathway meetings to increase surgical capacity at Royal Papworth. The slow pace of improvement means that substantial structural changes -- perhaps including expansion of surgical sessions, adoption of virtual MDTs, and streamlining of pre-operative investigations -- remain urgently required.

\section{Survival After IPAH Diagnosis Varies Dramatically by Sex}

The supplementary survival analysis examines Kaplan-Meier estimates for adults in the longitudinal cohort (first referral on or after 1 April 2009), stratified by sex, age group, and World Health Organization functional class. This cohort comprises 1,517 patients with idiopathic, heritable, or drug-induced pulmonary arterial hypertension, subdivided into those with and without comorbidities.

Among patients without comorbidities, women enjoy substantially better long-term survival than men. The five-year cumulative survival rate is 61\% for women versus only 42\% for men. By year six, female survival remains at 56\%, whereas male survival has fallen to 36\%. Women with comorbidities fare considerably worse than women without, at 40\% survival by year three, but they still outperform men of comparable complexity.

\begin{figure}[H]
  \centering
  \includegraphics[width=\textwidth]{./data/workspaces/17/chart_survival_sex.png}
  \caption{Kaplan-Meier survival curves from diagnosis stratified by sex and comorbidity status for IPAH patients. Solid lines denote IPAH without comorbidities; dashed lines denote IPAH with comorbidities. Female patients consistently exhibit higher cumulative survival rates across the first eight years. At five years post-diagnosis, women without comorbidities stand at 61\% survival while men drop to 42\%. Patients with comorbidities decline steeply regardless of sex, reaching 40\% or less by year three.}
\end{figure}

Demographic context helps explain part of the gap: men diagnosed with IPAH without comorbidities are older on average (median 65 years, mean 61) than women (median 57, mean 56), which partially accounts for the poorer outcomes. Nevertheless, the sex difference persists even when controlling for comorbidity presence, suggesting biological, therapeutic, or access-to-care factors warrant further investigation.

Patients with comorbidities face markedly worse prognoses. Their median time to 50\% mortality ranges from 2 years and 211 days (men) to 2 years and 244 days (women), compared with 3 years and 279 days for men without comorbidities and 7 years and 77 days for women without comorbidities. Comorbid patients tend to be older at diagnosis (median age 69--71) and present with more severe haemodynamics. The sample sizes for comorbid groups (n = 204 females, n = 181 males) are notably smaller than for those without comorbidities (n = 740 and n = 392 respectively), limiting the precision of estimates beyond three years.

\section{Age Is a Dominant Predictor of Long-Term Survival}

Within IPAH patients without comorbidities, age at diagnosis explains nearly as much variance in survival as comorbidity status itself. The cohort was divided into quintiles spanning ages 18 to 91, and five-year cumulative survival decreases monotonically from youngest to oldest:

\begin{table}[H]
\centering
\small
\caption{Five-year cumulative survival (%) by age quintile for IPAH without comorbidities}
\begin{tabularx}{\textwidth}{lrrr}
\toprule
Age Quintile & Age Range & Median Age at Diagnosis & Five-Year Survival (\%) \\
\midrule
Quintile 1 & 18--41 & 32 & 80 \\
Quintile 2 & 41--54 & 47 & 75 \\
Quintile 3 & 54--66 & 60 & 45 \\
Quintile 4 & 66--74 & 70 & 35 \\
Quintile 5 & 74--91 & 78 & 30 \\
\midrule
IPAH with comorbidities (all ages) & 26--90 & 70 & 25 \\
\bottomrule
\end{tabularx}
\end{table}

The survival gradient is steep: patients in the youngest quintile retain 80\% cumulative survival at five years, compared with only 30\% in the oldest quintile -- a gap of 50 percentage points. This pattern reflects both the biological impact of PH in elderly patients and the likelihood that older individuals harbour undiagnosed comorbidities that elevate pulmonary vascular resistance independently.

\begin{figure}[H]
  \centering
  \includegraphics[width=\textwidth]{./data/workspaces/17/chart_survival_age.png}
  \caption{Five-year survival trajectories by age at diagnosis for IPAH without comorbidities (Q1--Q5) compared with the combined group of IPAH with comorbidities. The youngest quintile (ages 18--41) maintains 80\% survival at year five, whereas the oldest quintile (74--91) drops to 30\%. Notably, the comorbidity group falls even lower at 25\%, demonstrating that the combination of advanced age and additional cardiovascular or respiratory disease creates the poorest prognosis profile.}
\end{figure}

Compared with any single age quintile among those without comorbidities, the comorbidities group performs worse than the two oldest quintiles alone, at just 25\% five-year survival. Given that comorbid patients' median age (70) falls squarely between quintiles 4 and 5, this result implies that comorbidity adds risk equivalent to advancing roughly one additional decade in age. Clinicians interpreting risk should therefore consider both chronological age and comorbidity burden jointly rather than treating either as sufficient alone.

\section{Clinical Severity at Presentation Determines Early Trajectories}

World Health Organization functional class provides another powerful prognostic signal. Class II patients are ambulatory without limitation of ordinary activity, class III have marked limitation, and class IV cannot carry out any physical activity without symptoms at rest.

For IPAH patients without comorbidities, the one-year survival difference between WHO class II and class IV is 10 percentage points (95\% versus 85\%), widening to a 27-point gap by year four (75\% versus 48\%). Patients presenting in class IV fare so poorly that survival data becomes sparse after two to three years, underscoring the urgency of earlier intervention.

\begin{figure}[H]
  \centering
  \includegraphics[width=\textwidth]{./data/workspaces/17/chart_survival_who_detail.png}
  \caption{Survival curves stratified by WHO functional class for IPAH without comorbidities. Class II patients maintain 75\% survival at year four and likely beyond (the plot truncates at five observations due to cohort boundaries). Class III patients reach 62\% at year four, falling to 43\% by year eight. Class IV patients drop to 48\% by year four, consistent with the poor prognosis expected for patients already classified as unable to perform any physical activity without dyspnoea or fatigue.}
\end{figure}

When comorbidities are present, the picture is even starker. Patients in WHO class III with comorbidities have only 50\% survival at three years, while those in class IV with comorbidities reach merely 55\% by two years. Small sample sizes prevent robust estimation beyond these early timepoints (the audit stops generating curves when the surviving cohort falls below 50 patients).

These patterns reinforce the importance of accurate initial functional classification and suggest that patients arriving in classes III or IV should receive aggressive combination therapy rather than sequential monotherapy trials, because delay increases early mortality risk.

\section{Centre-Level Variation Reveals Opportunities for Peer Learning}

While national aggregates tell one story, comparing individual centres reveals variation that can inform targeted quality improvement. Four centres stood out for specific performance profiles:

\begin{enumerate}[leftmargin=*]
\item \textbf{Newcastle} showed remarkable improvement on Standard 4 (seen within 30 days), rising from 19\% in 2016--17 to 78\% in 2018--19 -- an increase of 59 percentage points in two years. Their local action plan cites systematic monitoring of all new patient activity.
\item \textbf{Golden Jubilee} experienced a deterioration on Standard 8 (drug therapy within 12 weeks), falling from 84\% to 61\% over three years. They explicitly committed to addressing the backlog by increasing right heart catheterisation throughput.
\item \textbf{Imperial College} saw its cardiac catheterization rate (Standard 5) drop from 100\% to 92\%, falling below the 95\% target for the first time in the audit history. The local action plan notes that most excluded patients had undergone catheterization externally, and documentation lag caused apparent non-compliance. This highlights a data quality issue rather than a care quality problem.
\item \textbf{Royal Free} improved continuously on multiple fronts: Standard 3 (timely diagnosis) from 98\% to 99\%, Standard 8 (drug therapy) from 71\% to 87\%, and Standard 12 (annual consultation) from 96\% to 97\%. Their actions include quarterly standard reviews and identification of discrepancies closer to entry time.
\end{enumerate}

The most instructive comparison concerns Standard 4 (seen within 30 days), where Great Ormond Street (children's centre) achieved 100\% in 2018--19 and Golden Jubilee reached 74\%, while Royal Brompton languished at 57\% and Sheffield declined from 64\% to 55\%. The children's centre benefits from a dedicated administrative structure noted by Consultant Physician John Wort: ``We analyse our data against the standards every two months and identify where we need to improve. It is a real team effort.'' Adopting similar governance cadence could benefit adult centres currently lagging.

\section{Methods and Limitations}

\subsection{Data Sources and Construction}

This analysis draws primarily from the NAPH 10\textsuperscript{th} Annual Report (published 24 October 2019) and its Supplementary Survival Analysis. The main report contains 13 audit standards evaluated at national and centre level across audit years 2015--16 through 2018--19. The supplementary file provides Kaplan-Meier survival estimates stratified by sex, age quintile, and WHO functional class for the IDAH longitudinal cohort (first referral on or after 1 April 2009). Five local action plan documents from individual centres supplement the quantitative analysis with qualitative improvement commitments.

All tables and numerical values cited in this report were extracted directly from the published markdown files. Where exact figures required parsing from complex nested markdown table structures, manual verification against the source text confirmed alignment.

\subsection{Key Limitations}

\begin{enumerate}[leftmargin=*]
\item \textit{Rounding and aggregation.} National standard results are rounded to the nearest whole number before determining whether a target is met. This masking means that a centre reporting 89\% on an 80\% target is effectively indistinguishable from one at 84\% once rounded. Fine-grained comparisons between adjacent centres should therefore be treated cautiously.

\item \textit{Scottish patients and mortality tracing.} Scottish patients submitted by Golden Jubilee are pseudo-nymised in the NAPH database and cannot be traced at the Office for National Statistics or National Records of Scotland. Consequently, Scottish data are excluded from survival analyses entirely. This introduces potential geographic bias, because the UK-wide survival curves represent largely English and Welsh populations.

\item \textit{Survival curve truncation.} Curves terminate when the remaining cohort drops below 50 patients, producing abrupt endpoints that may underestimate true long-term survival for groups with small samples. The comorbidity strata are especially vulnerable to this effect, with most estimates ceasing at two to four years.

\item \textit{Observational design.} NAPH is an observational registry, not a randomised trial. Differences in survival between groups may reflect unmeasured confounders such as socioeconomic status, adherence to medication, or subtle variations in treatment intensity beyond what the audit captures.

\item \textit{Temporal stability of centre-specific findings.} Trends based on two or three years of data may reflect temporary staffing changes, IT migration issues, or transient referral patterns rather than persistent quality gaps. Repeated observation across future audit cycles is necessary before attributing causality to observed shifts.
\end{enumerate}

\subsection{Recommendations for Future Data Collection}

To strengthen confidence in these findings, the following additions to the NAPH dataset would be valuable:

\begin{itemize}[leftmargin=*]
\item Incorporation of Scottish patient outcomes through alternate mortality linkage agreements.
\item Recording of comorbidity scores (e.g., Charlson or Elixhauser indices) alongside binary comorbidity flags, enabling multivariable adjustment in survival modelling.
\item Collection of treatment escalation timelines to complement the binary drug-within-12-weeks metric, distinguishing monotherapy initiation from combination regimen optimisation.
\item Centralised tracking of pre-operative intervals from CTEPH diagnosis through PEA referral to final surgery date, enabling centre-level benchmarking against the 90\% standard.
\end{itemize}

\section{Summary of Key Findings}

\begin{table}[H]
\centering
\small
\caption{Condensed summary of principal findings with supporting figures}
\begin{tabularx}{\textwidth}{lccrX}
\toprule
Metric & Target & 2018--19 Result & Gap & Interpretation \\
\midrule
Timely Diagnosis (NS3) & $\geq$95\% & 99\% & $-$4pp & Excellent; 6 of 7 centres achieved full compliance \\
Seen Within 30 Days (NS4) & $\geq$50\% & 63\% & $-$13pp & Steady improvement; all centres met minimum threshold \\
Cardiac Cath (NS5) & $\geq$95\% & 97\% & $-$2pp & Strong; Imperial College slipped to 92\% \\
WHO/Exercise (NS6) & $\geq$90\% & 97\% & $-$3pp & Near-universal documentation excellence \\
Drug Within 12 Wks (NS8) & $\geq$80\% & 81\% & $-$1pp & Borderline; only 4 of 7 centres hit target \\
Diagnosis Recorded (NS9) & $\geq$99\% & 100\% & Met & Perfect compliance across all centres \\
PDE5 Inhibitor (NS10) & $\geq$80\% & 91\% & $-$11pp & Strong adherence to guideline-recommended first-line therapy \\
Quality of Life (NS11) & $\geq$90\% & 96\% & $-$6pp & Rapid uptake from 78\% in 2015--16 to 96\% in 2018--19 \\
Annual Consultation (NS12) & $\geq$90\% & 96\% & $-$6pp & Stable high performance maintained since 2015--16 \\
PEA Wait < 4 Months (NS13) & $\geq$90\% & 11\% & $-$79pp & Systemic bottleneck; requires surgical capacity expansion \\
\bottomrule
\end{tabularx}
\end{table}

\begin{table}[H]
\centering
\small
\caption{Survival at key milestones by patient subgroup}
\begin{tabularx}{\textwidth}{lcrX}
\toprule
Subgroup & One-Year Survival (\%) & Five-Year Survival (\%) & Sample Size (n) \\
\midrule
IPAH w/o comorb: Female & 88 & 61 & 740 \\
IPAH w/o comorb: Male & 85 & 42 & 392 \\
IPAH w/ comorb: Female & 80 & (stops at yr 3) & 204 \\
IPAH w/ comorb: Male & 75 & (stops at yr 3) & 181 \\
Q1 (age 18--41) & 98 & 80 & 227 \\
Q3 (age 54--66) & 90 & 45 & 226 \\
Q5 (age 74--91) & 80 & 30 & 226 \\
With comorbidities & 75 & 25 & 385 \\
\bottomrule
\end{tabularx}
\end{table}

These figures consolidate the central message: overall quality of care for pulmonary hypertension in Great Britain is high and steadily improving, but significant pockets of underperformance persist -- most critically in surgical access for CTEPH patients and in timely initiation of drug therapy for selected centres. Survival gradients by demographic and clinical characteristics underscore the need for risk-stratified management protocols.

\end{document}
